% ═════════════════════════════════════════════════════════════════════════════
\section{Methodology}
\label{sec:methodology}

\subsection{Hypotheses}

\begin{description}
  \item[H1] CS2 skin index returns are significantly correlated with Bitcoin
    returns (market integration hypothesis).

  \item[H2] Bitcoin Granger-causes CS2 skin prices, but not vice versa --
    motivated by the shared speculative, crypto-adjacent user demographic.

  \item[H3] The CS2 market responds to Bitcoin shocks with a lag of several
    periods (delay hypothesis).

  \item[H4] The CS2 market exhibits higher volatility than the S\&P 500 but
    similar volatility clustering patterns to Bitcoin.
\end{description}

\subsection{Data}

\subsubsection{CS2 Skin Price Index}
% TODO: Source: Steam Marketplace transaction data (Kieran, Kaggle, CC-BY-NC-SA-4.0).
% 22,495 items. Filtered to wearable items (weapon skins, knives, gloves).
% Liquidity filter: drop items with fewer than 2 trades per month in any
% calendar month.
% Aggregation: daily VWAP per item, then value-weighted index across all
% liquid items (weight = trading\_value / sum(trading\_value) per day).
% Result: one daily CS2 index level series covering 2013--2024.

\subsubsection{Comparison Assets}
% TODO: S\&P 500, Bitcoin (BTC-USD), Gold (GC=F) via \texttt{quantmod} /
% Yahoo Finance. Time period aligned to CS2 data. Frequency: daily.
% All series converted to log returns for analysis.

\subsection{Analytical Pipeline}

\subsubsection{Exploratory Data Analysis}
% TODO: Time series plots of index levels and log returns.
% Descriptive statistics (mean, SD, skewness, kurtosis).
% Correlation matrix. ACF/PACF plots.
% ADF test for stationarity; all analysis performed on log returns.

\subsubsection{GARCH(1,1) Volatility Modelling}
% TODO: Fit GARCH(1,1) to CS2, BTC, and S\&P 500 return series.
% Compare $\alpha + \beta$ (persistence) and unconditional volatility across
% assets to test H4. Extract conditional volatility series.

\subsubsection{VAR Model and Granger Causality}
% TODO: VAR($p$) on CS2, BTC, S\&P 500 log returns.
% Lag order selected via AIC/BIC.
% Granger causality tests: BTC $\to$ CS2, CS2 $\to$ BTC,
% S\&P 500 $\to$ CS2 (tests H1 and H2).

\subsubsection{Impulse Response Functions}
% TODO: Orthogonalised IRF from the estimated VAR.
% Response of CS2 returns to a 1-SD shock in BTC and S\&P 500.
% Time horizon 10--20 days; assess speed and persistence of absorption (H3).
